% ================= 1. INTRODUCTION =================
\section{Introduction}

\subsection{Background \& Problem Statement}
\label{sec:background}

The Cubli is a reaction-wheel inverted pendulum: a cube-shaped platform that balances on an edge or a corner using only internal actuation, with no external supports or contact forces beyond the single line or point of contact with the ground. The concept was introduced by the Institute for Dynamic Systems and Control (IDSC) at ETH Zurich \cite{gajamohan2012cubli, gajamohan2013cubli}, whose original prototype ($15\times15\times15$~cm) mounts one brushless DC motor and flywheel on each of three orthogonal faces. Commanding a change in wheel speed produces, by conservation of angular momentum, an equal and opposite reaction torque on the cube body; braking a wheel abruptly transfers its stored angular momentum to the body in a near-impulsive event, which is what allows the cube to \emph{jump up} from resting on a face into a balance on an edge or corner.

Dynamically, the system is an \emph{underactuated}, nonlinear, unstable-equilibrium control problem: three actuators (the wheels) must stabilize a body with more degrees of freedom than actuated inputs, about equilibria (edge and corner balance) that are open-loop unstable and increasingly coupled across axes as the contact geometry shrinks from a line (edge) to a point (corner). This is structurally the same problem as the classical inverted pendulum, extended from a single rotational axis to a fully three-dimensional, momentum-exchange-actuated body — which is why the Cubli has become a standard benchmark platform in the controls community over the past decade.

The problem this project addresses is therefore twofold: (i) derive and validate a dynamic model and controller capable of stabilizing this underactuated system about its unstable equilibria and driving it through commanded reorientations, and (ii) realize this in hardware as a modular, largely parametric design, since final actuator and sensor specifications are not yet confirmed at the time of writing. The target capability set, in order of increasing difficulty, is edge balance, corner balance, controlled rotation, and walking (Section~\ref{sec:objectives}).

\subsection{Purpose, Application \& Mission Context}
\label{sec:application}

The Cubli is not presented here as a balancing novelty but as a terrestrial, gravity-loaded demonstrator of \textbf{momentum-exchange attitude control} — the same physical mechanism that governs how real spacecraft point themselves. Nearly every three-axis-stabilized satellite holds and changes its attitude using reaction wheels: flywheels whose commanded speed change produces a reaction torque on the spacecraft body, exactly as in the Cubli. Spacecraft typically carry three or four such wheels (often in a pyramidal arrangement for redundancy) — directly analogous to the Cubli's three orthogonal wheels — and this reaction-wheel approach is the standard for CubeSats and small satellites in particular, making the Cubli representative of a real, currently-flown hardware class rather than an idealized toy problem. Reaction wheels are distinct from control moment gyros (CMGs), which gimbal a constantly-spinning wheel instead of changing its speed to produce much larger torques \cite{votel2012cmgvsrw}; positioning the Cubli explicitly as a \emph{reaction-wheel}, not CMG, demonstrator keeps the application claim precise.

Framed as an ADCS problem, the Cubli's core behaviors map directly onto a spacecraft's day-to-day pointing task:

\begin{itemize}
    \item \textbf{Edge / corner balance} $\rightarrow$ regulation to a commanded reference attitude and disturbance rejection — holding a pointing setpoint against perturbations.
    \item \textbf{Controlled rotation} $\rightarrow$ a slew maneuver: reference tracking to a new commanded attitude, exactly as a satellite re-aims a payload or antenna.
    \item \textbf{Jump / walk} $\rightarrow$ momentum-exchange hopping mobility, flight-proven on small bodies where wheeled traction is impossible. JAXA's Hayabusa2 mission deployed the MINERVA-II1 rovers \cite{yoshimitsu1999minerva} and the MASCOT lander \cite{ho2017mascot} on asteroid Ryugu (2018), both of which move by spinning up and braking an internal flywheel or eccentric mass to hop across the surface, with MASCOT additionally self-righting the same way. Walking, as a sequence of controlled, directional hops, is a genuine step beyond a single ballistic hop and is an active research frontier (e.g.\ ETH's SpaceHopper \cite{spiridonov2024spacehopper}).
\end{itemize}

This gives the project two complementary application narratives rather than one: it is simultaneously a low-cost, hardware-in-the-loop testbed for standard spacecraft \emph{attitude determination and control} (ADCS), and a ground analog for the momentum-exchange \emph{surface mobility} now used on small-body missions. Framed against a companion star-tracker project, which answers \emph{``where am I pointing?''} (attitude determination), the Cubli answers \emph{``drive me to that reference and hold it''} (attitude control) — the two halves of a complete ADCS loop.

\subsection{Objectives}
\label{sec:objectives}

The project targets four behaviors, deliberately ordered as a difficulty ramp in which each milestone is a control precondition for the next:

\begin{enumerate}
    \item \textbf{Edge balance} — stabilize the cube on a single edge, the least unstable equilibrium (contact is a line, one primary axis of instability).
    \item \textbf{Corner balance} — stabilize the cube on a single vertex, the fully unstable equilibrium (point contact, all three axes coupled and unstable simultaneously).
    \item \textbf{Controlled rotation} — execute a commanded slew between equilibria (e.g.\ edge to edge, or in-place reorientation), demonstrating reference tracking rather than pure regulation.
    \item \textbf{Walking} — chain repeated, directional controlled falls between edges/corners into a net translation across a surface, re-stabilizing after each transition.
\end{enumerate}

These map one-to-one onto the applications of Section~\ref{sec:application}: (1)–(2) demonstrate pointing regulation and disturbance rejection of increasing difficulty; (3) demonstrates slew/reference-tracking; (4) demonstrates momentum-exchange mobility of the kind flown on Hayabusa2. Milestones (1)–(3) are treated as committed deliverables for this design; walking (4) is scoped as the stretch objective, contingent on hardware confirmation and the jump-up feasibility check (angular momentum required to tip the cube vs.\ available motor/wheel torque), to be resolved once final components are known.